\chapter{Źródła}

\begin{lstlisting}[caption={Wybrane metody kontekstu kompilatora}, label={lst:ok-wybrane}, morekeywords={type, string, interface, any, map, func}]
func (st *ScopeTree) GetScopeByName(name string) *ScopeTree {
	for _, c := range st.GetChildren() {
		if c.name == name {
			return c
		}
		foundScope := c.GetScopeByName(name)
		if foundScope != nil {
			return foundScope
		}
	}
	return nil
}

func SearchIdUp(scope *ScopeTree, id string) *ScopeTree {
	if scope == nil {
		return nil
	}
	vars := *scope.GetVars()
	if _, found := vars[id]; found {
		return scope
	}
	return SearchIdUp(scope.GetParent(), id)
}

func GetWhatFunc(scope *ScopeTree) *ScopeTree {
	if scope.scopeType == FUNC {
		return scope
	}
	return GetWhatFunc(scope.GetParent())
}

func GetWhatFor(scope *ScopeTree) *ScopeTree {
	if scope.scopeType == FOR {
		return scope
	}
	return GetWhatFor(scope.GetParent())
}

func BlockScopeName(name string) string {
	scopeName := fmt.Sprintf("BLOCK__SCOPE__%s__%d", name, scopeId)
	scopeId++
	return scopeName
}

func EndFnLabel(name string) string {
	return fmt.Sprintf("END__%s", name)
}

\end{lstlisting}

\begin{lstlisting}[caption={Metoda Visit drzewa składniowego}, label={lst:ok-visit}, morekeywords={type, string, interface, any, map, func}]
func (v *Visitor) Visit(tree antlr.ParseTree) any {
	switch ctx := tree.(type) {
	case *parsing.ProgramContext:
		return v.VisitProgram(ctx)
	case *parsing.FuncContext:
		return v.VisitFunc(ctx)
	case *parsing.ArgsContext:
		return v.VisitArgs(ctx)
	case *parsing.StmtContext:
		return v.VisitStmt(ctx)
	case *parsing.ExprContext:
		return v.VisitExpr(ctx)
	case *parsing.ParamContext:
		return v.VisitParam(ctx)
	case *parsing.Effect_blockContext:
		return v.VisitEffect_block(ctx)
	case *parsing.BlockContext:
		return v.VisitBlock(ctx)
	case *parsing.For_stmtContext:
		return v.VisitFor_stmt(ctx)
	case *parsing.If_stmtContext:
		return v.VisitIf_stmt(ctx)
	case *parsing.Func_callContext:
		return v.VisitFunc_call(ctx)
	case *parsing.Call_argsContext:
		return v.VisitCall_args(ctx)
	default:
		return ctx.Accept(v)
	}
}
\end{lstlisting}

\begin{lstlisting}[caption={Wybrana metoda wizytująca}, label={lst:ok-wybrane-wizytor}, morekeywords={type, string, interface, any, map, func}]
func (v *Visitor) VisitFor_stmt(ctx *parsing.For_stmtContext) any {
	scopeName := BlockScopeName("for")
	startLabel := StartBlockLabel(scopeName)
	endLabel := EndBlockLabel(scopeName)
	stmtLabel := EndStmtLabel(scopeName)

	parent := v.cctx.GetCurrScope()
	scope := NewScopeTree(scopeName, parent, FOR)
	v.cctx.currentScope = scope

	v.Visit(ctx.Stmt(0))
	v.GenLabel(startLabel)
	// init variable
	// eval expr
	condition := v.Visit(ctx.Expr()).(*ScopeVar)
	// check condition, break
	v.GenCmpAddrReg(condition.offset, "0")
	v.GenJz(endLabel)

	// visit stmts
	v.Visit(ctx.Block())

	// increment
	v.GenLabel(stmtLabel)
	v.Visit(ctx.Stmt(1))

	// jump to start
	v.GenJmp(startLabel)
	v.GenLabel(endLabel)

	v.cctx.currentScope = parent
	return nil
}
\end{lstlisting}

\begin{lstlisting}[caption={Struktury generacji kodu}, label={lst:gen-struktury}, morekeywords={type, string}]
type Instruction struct {
	op string
	dst *string
	src *string
}

type Codegen struct {
	global []Instruction
	text []Instruction
	data []Instruction
};
\end{lstlisting}

\begin{lstlisting}[caption={Generacja kodu asemblera z instrukcji}, label={lst:gen-generacja}, morekeywords={func, string}]
func instrArrToString(arr []Instruction) string {
	asm := ""
	for _, instr := range arr {
		asm = fmt.Sprintf("%s\n%s", asm, instr.toString())
	}
	return asm
}
\end{lstlisting}

\begin{lstlisting}[caption={Interfejs wiersza poleceń}, label={lst:cli}, morekeywords={type, func, string}]
type Options struct {
	Input string
	Output string
	IncludeDir string
}

func NewRootCmd() (*cobra.Command, *Options) {
	opts := &Options{}
	cmd := &cobra.Command {
		Use: "blunderbuss",
		Short: "A compiler for blunderbuss",
	}

	cmd.Flags().StringVarP(&opts.Input, "input", "i", "", "source file")
	cmd.Flags().StringVarP(&opts.Output, "output", "o", "a.asm", "output file")
	cmd.Flags().StringVarP(&opts.IncludeDir, "include", "I", "", "include directory")

	cmd.MarkFlagRequired("input")
	cmd.MarkFlagRequired("include")

	return cmd, opts
}
\end{lstlisting}

\begin{lstlisting}[caption={Warstwa pośrednia dla preprocesowania GCC}, label={lst:cli-gcc}, morekeywords={func, string, nil}]
func Preprocess(filename string, includeDir string) []byte {
	fmt.Println("[BBUSS] preprocess and compiler step...")

	cmd := exec.Command("gcc", "-E", "-P", "-x", "c", "-I"+includeDir, filename)
	out, err := cmd.CombinedOutput()
	if err != nil {
		fmt.Println("[ERR] failed to preprocess with GCC")
		fmt.Println(string(out))
		os.Exit(1)
		return nil
	}
	fmt.Println(string(out))

	return out
}
\end{lstlisting}

\begin{lstlisting}[caption={Definicja błędów parsera}, label={lst:err-parser}, morekeywords={type, string, nil}]
// error.go
type SyntaxError struct {
    Msg  string
    Text string
    Line int
    Col  int
}

type CustomErrorListener struct {
    *antlr.DefaultErrorListener
    Errors []SyntaxError
}

func (l *CustomErrorListener) SyntaxError(
    recognizer antlr.Recognizer,
    offendingSymbol interface{},
    line, column int,
    msg string,
    e antlr.RecognitionException,
) {
    tokenText := ""
    if t, ok := offendingSymbol.(antlr.Token); ok {
        tokenText = t.GetText()
    }

    l.Errors = append(l.Errors, SyntaxError{
        Msg:  msg,
        Text: tokenText,
        Line: line,
        Col:  column,
    })
}

// main.go
parser.RemoveErrorListeners()
listener := &semantics.CustomErrorListener{}
parser.AddErrorListener(listener)

if len(listener.Errors) > 0 {
	fmt.Println("[ERR] Syntax errors found")
	for _, e := range listener.Errors {
		fmt.Printf("Syntax error at %d:%d near '%s': %s\n", 
			e.Line, e.Col, e.Text, e.Msg)
	}
	os.Exit(1)
}

\end{lstlisting}

\begin{lstlisting}[caption={Definicja błędów kompilatora}, label={lst:err-compiler}, morekeywords={type, string, error, nil, func}]
//error.go
type CompilerError interface {
	error
	Context() antlr.ParserRuleContext
}

type BaseCompilerError struct {
	Msg string
	Ctx antlr.ParserRuleContext
}

func (e *BaseCompilerError) Error() string {
	token := e.Ctx.GetStart()
	return fmt.Sprintf("Compiler error at line %d:%d: %s", 
		token.GetLine(), token.GetColumn(), e.Msg)
}

func (e *BaseCompilerError) Context() antlr.ParserRuleContext {
	return e.Ctx
}

type UndefinedVariableError struct {
	*BaseCompilerError
	VarName string
}

type InitializedVariableError struct {
	*BaseCompilerError
	Type string
	VarName string
}

type TypeMismatchError struct {
    *BaseCompilerError
    Expected string
    Got string
}

type IndexError struct {
	*BaseCompilerError
	Type string
	VarName string
}

type FnMismatchError struct {
	*BaseCompilerError
	FnName string
	Expected string
	Got string
	Reason string
}

//main.go
visitor := semantics.NewBlunderbussVisitor()
if len(visitor.Errors) > 0 {
	fmt.Println("[ERR] Compilation failed with errors:")
	for _, e := range visitor.Errors {
		fmt.Println(e)
	}
	os.Exit(1)
}
\end{lstlisting}

\begin{lstlisting}[caption={Memoizacja w standardowej bibliotece}, label={lst:sb-mem}, morekeywords={func, nil, string, C.long, int64, var, any}]
//export ____SetM
func ____SetM(values *C.long, types *C.long, length C.long) {
	var goArgs []any
	var i int64
	cptr := C.malloc(8)
	for i = 0; i < int64(length); i ++ {
		val := *(*C.long)(unsafe.Pointer(uintptr(unsafe.Pointer(values)) + uintptr(i)*unsafe.Sizeof(*values)))
		t := *(*C.int)(unsafe.Pointer(uintptr(unsafe.Pointer(types)) + uintptr(i)*unsafe.Sizeof(*types)))

		switch t {
		case 1:
			if i == int64(length) - 1 {
				*(*int64)(cptr) = int64(val)
			}
			goArgs = append(goArgs, int64(val))
		case 2:
			str := C.GoString((*C.char)(unsafe.Pointer(uintptr(val))))
			goArgs = append(goArgs, str)
			if i == int64(length) - 1 {
				cstr := C.CString(str)
				*(*uintptr)(cptr) = uintptr(unsafe.Pointer(cstr))
			}
		case 3:
			ptr := unsafe.Pointer(uintptr(val))
			goArgs = append(goArgs, ptr)
		case 0:
			ptr := unsafe.Pointer(uintptr(val))
			goArgs = append(goArgs, ptr)
		default:
		}
	}
	var key uint64 
	if len(goArgs) == 1 {
		key = Hash(goArgs[0].(unsafe.Pointer))
	} else {
		key = Hash(goArgs[0].(unsafe.Pointer), goArgs[1:len(goArgs)-1]...)
	}


	m[key] = cptr
}

//export ____GetM
func ____GetM(values *C.long, types *C.long, length C.long) unsafe.Pointer {
	var goArgs []any
	var i int64
	for i = 0; i < int64(length); i ++ {
		val := *(*C.long)(unsafe.Pointer(uintptr(unsafe.Pointer(values)) + uintptr(i)*unsafe.Sizeof(*values)))
		t := *(*C.int)(unsafe.Pointer(uintptr(unsafe.Pointer(types)) + uintptr(i)*unsafe.Sizeof(*types)))

		switch t {
		case 1:
			goArgs = append(goArgs, int64(val))
		case 2:
			str := C.GoString((*C.char)(unsafe.Pointer(uintptr(val))))
			goArgs = append(goArgs, str)
		case 3:
			ptr := unsafe.Pointer(uintptr(val))
			goArgs = append(goArgs, ptr)
		case 0:
			ptr := unsafe.Pointer(uintptr(val))
			goArgs = append(goArgs, ptr)
		default:
		}
	}

	var key uint64 
	if len(goArgs) == 1 {
		key = Hash(goArgs[0].(unsafe.Pointer))
	} else {
		key = Hash(goArgs[0].(unsafe.Pointer), goArgs[1:len(goArgs)-1]...)
	}

	if val, ok := m[key]; ok {
		return val
	}
	return nil

}
\end{lstlisting}

\begin{lstlisting}[caption={Inne definicje funkcji w standardowej bibliotece}, label={lst:sb-inne}, morekeywords={func, C.long, int64, int, string, nil}]
//export Itoa
func Itoa(num C.long) *C.char {
	s := strconv.FormatInt(int64(num), 10)
	return C.CString(s)
}

//export PrintlnStr
func PrintlnStr(arg *C.char) {
	fmt.Println(C.GoString(arg))
}

//export PrintlnInt
func PrintlnInt(arg C.long) {
	fmt.Println(arg)
}
\end{lstlisting}
